%=======================================================================
% Chapter 5 -- draft (rewritten 2026-10-02).
%
% Scope (decided 2026-10-01/02): growth for ONE and TWO species, coupled into
% the existing IKM element following Klempt et al. 2024; species composition
% from the point model of Klempt et al. 2026. More species and all Abaqus
% work are the Keio continuation and do NOT appear here as results.
% Writing rules: CLAUDE.md, "Slides, notes and other documents" (first person
% singular, no Abaqus as thesis work, units on every constant, no internal
% run labels).
%
% HOW TO USE THIS FILE
%   \TierA blocks state results that are final. \TierB blocks hold numbers
%   that must be re-run before submission; the comment above each says why.
%   Delete both \newcommand definitions once the chapter is final.
%
% PORTING (see PORTING.md). Material to merge into the thesis tree:
%     - packages: amsmath, amssymb, bm, xcolor, graphicx. Nothing else.
%     - every label is namespaced ch5-*.
%     - one path to set: \ansysfigdir below.
%
% CITATIONS
%   growth field Eq. 34-36, kinematics, Table 2   -> Klempt2024DiffusionDrivenGrowth
%   multi-species point model, Table 1 cases       -> Klempt2026ContinuumBacterialGrowth
%   same point model, Bayesian updating            -> Fritsch2025BayesianMicrofilms
%   The 2024 and 2026 papers are different papers; do not swap them.
%=======================================================================
% ---- the one thing to set when porting -------------------------------
\providecommand{\ansysfigdir}{../assets/}
% ----------------------------------------------------------------------

% Drafting markers. Delete both once the chapter is final.
\providecommand{\TierA}[1]{\marginpar{\footnotesize A}#1}   % final
\providecommand{\TierB}[1]{\marginpar{\footnotesize B}{\color{gray}#1}} % to re-run

\chapter{Coupling biofilm growth into an existing FEM element}
\label{ch:ansys}

%-----------------------------------------------------------------------
\section{Scope and contribution}
\label{sec:ch5-scope}

% TierA. THESIS_ASSIGNMENT.md §1; scope decision 2026-10-01. Shortened 3 Oct
% 2026: the objectives are in chapter 1, the model equations in chapter 2.
\TierA{%
This chapter brings the point model of Chapter~\ref{ch:heine} into the
biofilm element developed at IKM for the partner project, which solves the
field model of Klempt et al.~\cite{Klempt2024DiffusionDrivenGrowth}
(Section~\ref{sec:theory_field}) in ANSYS. It covers one and two species, for
which the papers give numbers to check against, and contains the second and
third contribution of the thesis (Section~\ref{sec:intro_structure}):
\begin{itemize}
\item a growth law, $\mathbf F=\mathbf F_e\mathbf F_g$ with
      $\mathbf F_g=\alpha\mathbf I$, implemented as an ANSYS user material
      and verified against closed forms (\S\ref{sec:ch5-model}--\ref{sec:ch5-verification}),
      and coupled to the element's biofilm field through the growth equation
      at every Gauss point (\S\ref{sec:ch5-coupling});
\item for two species, the composition from the point model of Klempt et
      al.~\cite{Klempt2026ContinuumBacterialGrowth}
      (Section~\ref{sec:theory_point}), run at every Gauss point alongside the
      field (\S\ref{sec:ch5-composition}).
\end{itemize}}

%-----------------------------------------------------------------------
\section{The constitutive model}
\label{sec:ch5-model}

% TierA. Kinematics: Klempt2024DiffusionDrivenGrowth.
\TierA{%
Growth enters multiplicatively,
$\mathbf F = \mathbf F_e \mathbf F_v \mathbf F_g$ with
$\mathbf F_g = \alpha\mathbf I$ and $\alpha(0)=1$, as in Klempt et
al.~\cite{Klempt2024DiffusionDrivenGrowth}; $\alpha-1$ is the growth, and
$\alpha=1$ means no growth. The elastic response is compressible
Mooney--Rivlin, reducing to neo-Hookean at $C_{01}=0$. A viscous branch, a
deviatoric Newtonian dashpot carrying $\mathbf F_v$, is implemented but
switched off ($\eta=0$) in every result of this chapter, as in the paper.}

\TierA{%
Written out in the order the implementation evaluates it, a step takes the
total deformation gradient $\mathbf F$ and the viscous state
$\mathbf F_v^{\,n}$ and returns the Cauchy stress and $\mathbf F_v^{\,n+1}$:
\begin{enumerate}
\item the mechanical part removes the growth,
      $\mathbf F_{\mathrm{mech}} = \mathbf F\,\mathbf F_g^{-1}$, and a trial
      elastic gradient follows from the viscous state at the start of the step,
      $\mathbf F_e^{\mathrm{tr}} = \mathbf F_{\mathrm{mech}}(\mathbf F_v^{\,n})^{-1}$;
\item with $\mathbf b_e = \mathbf F_e \mathbf F_e^{\mathsf T}$ and
      $J_e = \det \mathbf F_e$, the deviatoric Kirchhoff stress follows from
      $W = C_{10}(\bar I_1 - 3) + C_{01}(\bar I_2 - 3) + \tfrac{1}{D_1}(J_e-1)^2$;
\item the viscous state advances by one step of the deviatoric flow rule,
      \begin{equation}
        \mathbf F_v^{\,n+1} =
        \Bigl( \mathbf I + \frac{\Delta t}{2\,\eta\,J_e}\,
        \bm\tau_{\mathrm{dev}} \Bigr)\mathbf F_v^{\,n},
        \label{eq:ch5-flow}
      \end{equation}
      with $\bm\tau_{\mathrm{dev}}$ taken at the trial state;
\item the stress is recomputed at $\mathbf F_v^{\,n+1}$, the volumetric term
      $\tfrac{2}{D_1}(J_e-1)$ is added, and the result is returned in the
      solver's Voigt ordering.
\end{enumerate}

\begin{figure}[htbp]
  \centering
  \includegraphics[width=0.9\linewidth]{\ansysfigdir flow_growth_kinematics.png}
  \caption[Multiplicative growth kinematics]{%
    The configurations the decomposition introduces. Growth carries the
    reference state to an intermediate one that is stress free but generally
    incompatible; the elastic part restores compatibility, and this is where
    the residual stress comes from.}
  \label{fig:ch5-kinematics}
\end{figure}

A body that can grow freely therefore stays stress free; stress appears only
where growth is constrained, by boundary conditions or by neighbouring
material that grows less.}

% TierA. Klempt 2024 Eq. 20 and Table 2; biofilm_material_v01.f, sYoungL < 0.
\TierA{%
\paragraph{Stiffness.} Klempt et al.\ weight the elastic energy with
$\phi^2$, so that the void, $\phi=0$, has no stiffness, and give
$\mu=3.3557$~Pa ($E=10$~Pa, $\nu=0.49$) in their Table~2. The routine
supports this form,
\begin{equation}
  E(\phi) = (\phi^2+f)\,E_{\mathrm{bio}},\qquad \nu=\nu_{\mathrm{bio}},
  \label{eq:ch5-stiffness}
\end{equation}
where $f$ is a small floor that keeps the void solvable; it is not in the
paper, whose void has no stiffness. It is tested to
reproduce the plain law fed the same $E$ at every $\phi$ and to scale the
stress by $\phi^2$. The paper's material is incompressible, with $\det\mathbf F_e=1$ enforced by
a Lagrange multiplier, and $\nu=0.49$ enters there only through $\mu$; here
the constraint is approximated by $\nu=0.49$. The element, SOLID185 in its default
$\bar{\mathbf B}$ formulation, uses one dilatation per element, the same idea as
the H1P0 element of the paper. A Python solution of the same problem shows
that with this formulation the mean stress changes by less than $10\,\%$
between $\nu=0.49$ and $0.499$, while full integration would overestimate it
several times (\S\ref{sec:ch5-results}).}

%-----------------------------------------------------------------------
\section{Implementation}
\label{sec:ch5-impl}

% TierA. ansys_usermat/usermat_biofilm.f, biofilm_material_v01.f.
\TierA{%
The law is implemented as an ANSYS \texttt{USERMAT}.%
\footnote{\texttt{ansys\_usermat/usermat\_biofilm.f} and
\texttt{biofilm\_material\_v01.f} in the accompanying repository.}
It is delivered to the partner framework as a routine with the same argument
shape as the framework's own elastic law: per-phase $(E,\nu)$ blended by the
local biofilm fraction, returning Cauchy stress and consistent tangent. The
delivered routine is an adapter around the verified core, and this is tested
rather than assumed: fed $(E,\nu)$, it reproduces the core fed the
$(C_{10},C_{01},D_1)$ those map to, at zero tolerance. It uses no solver
includes and no common blocks, so it does not depend on the ANSYS release.}

% TierA. tests/test_tangent_quality.py -- 648 combinations vs exact AD.
\TierA{%
The consistent tangent is obtained by perturbing $\mathbf F$ rather than in
closed form. A tangent that is slightly wrong does not give a wrong answer but
a Newton iteration that converges slowly or not at all, which is hard to
diagnose inside someone else's framework. I therefore compared it with
forward-mode automatic differentiation over the range the routine is expected
to see: stiffness over three decades, Poisson ratios up to $0.49$, growth up to
$\alpha-1=0.35$, both material paths, elastic and viscous. Over 648 combinations
the relative difference has median $1.2\times10^{-7}$ and worst case
$1.1\times10^{-6}$, with no part of the range degrading.}

%-----------------------------------------------------------------------
\section{Verification of the growth law}
\label{sec:ch5-verification}

\begin{figure}[htbp]
  \centering
  \includegraphics[width=0.95\linewidth]{\ansysfigdir flow_vv_thesis.png}
  \caption[Where this chapter's checks sit]{%
    The verification and validation hierarchy with the evidence of this
    chapter placed on it. All of it is verification, that is, whether the
    implementation solves the equations it states. Validation against a
    measured stress is not attempted.}
  \label{fig:ch5-vv}
\end{figure}

\subsection{Growth kinematics against closed forms}
% TierA. ansys_usermat/apdl/ free and constrained decks; closed_form_reference.py.
\TierA{%
Two growth cases have closed forms. Traction-free isotropic growth expands the
body to $\mathbf F=\alpha\mathbf I$ with zero stress. Fully constrained
growth gives an isotropic elastic deformation, so the deviatoric stress
vanishes and only
\begin{equation}
  \bm\sigma = \frac{2}{D_1}\,(J_e-1)\,\mathbf I,\qquad J_e=\alpha^{-3}
  \label{eq:ch5-constrained}
\end{equation}
remains, independently of the shear constants and of the viscosity.

In ANSYS 2022~R2, one SOLID185 element at $\alpha=1.05$ returns a largest
stress component of $1.9\times10^{-11}$ under free growth, seven decades below
the constrained answer. Under full constraint it returns
$\sigma_{xx}=\sigma_{yy}=\sigma_{zz}=-1.0193\times10^{-4}$, with shear
components at machine noise. The free case is not redundant: full constraint
can hide a sign error in $\mathbf F_e=\mathbf F\mathbf F_g^{-1}$, because the
element cannot move, and only the free case exposes it.}

\TierA{%
The constrained value is not equation~\eqref{eq:ch5-constrained} alone. It
equals that term plus a spherical term that I predicted in closed form from
the implementation (\S\ref{sec:ch5-verification-limits}), to every digit
ANSYS prints, with the extra term making up $47\,\%$ of the reported stress
for this material (Figure~\ref{fig:ch5-closedform}).

\begin{figure}[htbp]
  \centering
  \includegraphics[width=\linewidth]{\ansysfigdir v222_closed_form.png}
  \caption[The two growth cases, solved in ANSYS]{%
    One SOLID185 element, $\alpha=1.05$, the same material, differing only
    in the boundary conditions. Left: every stress component of both solves.
    Right: the constrained value decomposed into the continuum term and the
    spherical term of \S\ref{sec:ch5-verification-limits}. Only their sum
    reproduces the solver. The spherical term does not enter the von Mises
    stress.}
  \label{fig:ch5-closedform}
\end{figure}}

\subsection{The delivered routine on a larger mesh}
% TierA. Curved two-layer shell, 12 240 elements, adapter vs core in ANSYS.
\TierA{%
A two-layer curved shell of $12\,240$ elements, a non-growing substrate under
a growth layer, was solved twice in ANSYS: once with the verified core and once
with the delivered routine. The von Mises fields agree to every digit reported
(minimum $4.3523\times10^{-9}$, maximum $1.1862\times10^{-5}$, mean
$9.14434842156836\times10^{-6}$), with no solver errors. This is a consistency
check between two builds, not a comparison with a closed form.

The same model showed a meshing error that I corrected. The two layers had been
meshed separately and joined by merging coincident nodes, and on the curved
sweeps only $63$ of $5\,673$ interface nodes coincided, three axial lines. The
growth layer was held along those lines and free between them; the ``two-lobe''
bulge an earlier draft reported as a buckling mode was the layer lifting
between them. With one shared interface mesh the bulge is near uniform and
converges at second order: the $\theta$-profile error falls from $2.7\,\%$ to
$0.58\,\%$ over $2.4$k, $19$k and $150$k elements, and every growth value up
to $\alpha-1=0.2$ converges (Figure~\ref{fig:ch5-alpha-sweep}).

\begin{figure}[htbp]
  \centering
  \includegraphics[width=\linewidth]{\ansysfigdir growth_cylinder_alpha_sweep.png}
  \caption[The bonded curved shell under increasing growth]{%
    The bonded two-layer shell ($2.4$k elements) under prescribed growth
    $\alpha-1 = 0.01$ to $0.2$. Left: maximum radial displacement against
    $\alpha-1$, with the straight line through the $\alpha-1=0.01$ run. Right: the
    radial displacement across the arc divided by $\alpha-1$; the curves
    coincide up to $\alpha-1\approx0.05$.}
  \label{fig:ch5-alpha-sweep}
\end{figure}

The finite-difference tangent holds up inside the solver as well
(Figure~\ref{fig:ch5-convergence}): every substep converges in one to four
equilibrium iterations, and the automatic time stepping enlarges the increment
at every opportunity, from $0.5$ to $1.6875$, without a single cut-back.

\begin{figure}[htbp]
  \centering
  \includegraphics[width=\linewidth]{\ansysfigdir v222_convergence.png}
  \caption[Newton convergence with the finite-difference tangent]{%
    Convergence of the $12\,240$-element solve, read from the solver's
    listing. Left: force residual against equilibrium iteration, one curve per
    substep, with the ANSYS criterion dashed. Right: the time increment chosen
    for each substep and the iterations it needed.}
  \label{fig:ch5-convergence}
\end{figure}}

\subsection{What these checks establish, and what they do not}
\label{sec:ch5-verification-limits}
% TierA. DEVIATOR_SCALING_FINDING.md.
\TierA{%
The isochoric split in the implementation applies $J^{-2/3}$ to the subtracted
trace but not to the tensor beside it. The coded and the correct forms agree
at $J=1$ and nowhere else, and growth moves $J_e$ away from $1$ by
construction. The difference is purely spherical for any deformation, so it is
a pressure error: the deviatoric stress, and with it the von Mises stress
reported in this work, is unaffected at a given deformation. It is the
spherical term of Figure~\ref{fig:ch5-closedform}.

Two of the usual checks cannot see such an error. The free-growth case lands
exactly on $J_e=1$, where both forms agree, and reference values generated by
calling the implementation record its behaviour rather than test it. Only the
comparison with an independently derived closed form finds it, which is why
equation~\eqref{eq:ch5-constrained} is derived from the continuum statement
and not from the code. A wrong pressure can still reach the von Mises stress
indirectly, through the displacement field of an equilibrium solve; this route
is not measured here. It shrinks as $\nu\to0.5$: at the paper's $\nu=0.49$
the extra term is $1.3\,\%$ of the constrained pressure, against $31\,\%$ at
$\nu=0.3$.}

%-----------------------------------------------------------------------
\section{Coupling to the element's biofilm field}
\label{sec:ch5-coupling}

\subsection{The framework and the growth equation}
% TierA. OLIVER_MODEL_NOTES.md; Klempt 2024 Eq. 34-36.
\TierA{%
The partner framework solves the mechanics in ANSYS and its transport fields in
its own routines, alternating substep by substep
(Figure~\ref{fig:partner-loop}). The biofilm field follows Klempt et
al.~\cite{Klempt2024DiffusionDrivenGrowth}, Eq.~34 and 36:
\begin{equation}
  \dot\phi = \beta\,\Delta\phi + k_\alpha\,\alpha
   - r\,\frac{c}{k+c}\,\nabla\phi\cdot\frac{\nabla c}{|\nabla c|},
  \qquad \dot\alpha = k_\alpha\,\phi ,
  \label{eq:ch5-field}
\end{equation}
with diffusion, a local source, and a front term that moves the biofilm
boundary up the nutrient gradient. The element computes $\phi$ at its
integration points; I leave it unchanged. Its Laplacian and the gradients of
$\phi$ and $c$ are approximated from the values at the neighbouring
integration points, by a second-order Taylor expansion fitted in the weighted
least-squares sense, the Neighbored Element Method
\cite{Blaszczyk2022TaylorWLS,Rudolf2025NEM}. The field is updated
explicitly, which limits the time step to $h^2/(2\Delta t)\ge\beta$
\cite{Rudolf2025NEM}; with $\beta=0.02$~mm$^2/T^*$ this holds for every mesh
and time step used in this chapter.

\begin{figure}[htbp]
  \centering
  \includegraphics[width=\linewidth]{\ansysfigdir flow_oliver_solution_loop.png}
  \caption[Partner framework solution loop]{%
    Solution loop of the partner framework, traced from its source, and the
    point at which the growth law of this chapter attaches.}
  \label{fig:partner-loop}
\end{figure}

What I add sits in the material routine, at every Gauss point: take $\phi$
from the field, integrate the second equation of~\eqref{eq:ch5-field},
\begin{equation}
  \alpha_{n+1}=\alpha_n+k_\alpha\,\phi\,\Delta t ,
  \label{eq:ch5-eq36}
\end{equation}
and pass $\alpha$ to the growth law. I use the paper's $k_\alpha=10^{-3}$ per unit
of normalised time $T^*$, in the field and in~\eqref{eq:ch5-eq36} alike, and
the deck time is $T^*$.

The paper prints a second form of this law (Section~\ref{sec:theory_field}). Table~1, step~3a, of Klempt et
al.\ 2024 integrates $\alpha$ from the change of $\phi$,
$(\alpha_{n+1}-\alpha_n)/((1+\alpha_{n+1})\Delta t)=k_\alpha(\phi_{n+1}-\phi_n)/\Delta t$
(with $k_\alpha/\eta_\alpha$ read as $k_\alpha$), while its text states that
$\alpha$ follows Eq.~36. The two differ in where growth happens. On the
reproduced field of test case~4.1, at $T^*=1$, biofilm present from the start
reaches $\alpha-1=9.9\times10^{-4}$ with Eq.~36 and stays at zero with
Table~1, and newly colonised biofilm reaches $5.4\times10^{-4}$ and
$8.5\times10^{-4}$ (Figure~\ref{fig:ch5-alpha-law}). A seed that is full
from the start, as in the runs of this chapter, grows only under Eq.~36. I
follow the text, Eq.~36, as does the partner's element.

\begin{figure}[htbp]
  \centering
  \includegraphics[width=\linewidth]{\ansysfigdir fig_klempt2024_alpha_law.png}
  \caption[Eq.~36 against Table~1 of Klempt et al.\ 2024]{%
    The two growth laws of Klempt et al.\ 2024 on the reproduced field of
    test case~4.1 (Python). Left and middle: $\alpha-1$ at $T^*=1$ on the
    mid-plane. Right: mean over biofilm present from the start and over
    newly colonised biofilm.}
  \label{fig:ch5-alpha-law}
\end{figure}}

\subsection{Growth counted once per increment}
% TierA. phi_trace checks on IKMHIWI03; 46 calls / 11 substeps.
\TierA{%
ANSYS calls the material routine several times per substep, once per
equilibrium iteration and once more for output, and hands back the converged
state each time. Equation~\eqref{eq:ch5-eq36} must therefore start from the
stored state variable and must not accumulate over iterations. A trace of every
call (46 calls over 11 substeps in the first run) confirms that growth is
counted once per increment, that each increment starts from the previous one's
end value, and that $\alpha$ follows~\eqref{eq:ch5-eq36} exactly.}

\subsection{Four properties of the framework found along the way}
% TierA. ONE_SPECIES_COUPLING.md, stage 3/4 and result-set-1 sections.
\TierA{%
\begin{enumerate}
\item \emph{The element's own growth variable is a mean.} The element already
      integrates $\dot\alpha=k_\alpha\phi$ for each of its two species fields
      and uses their mean as the growth of the material. With one species the
      unused field stays at $\alpha=1$, so the growth is halved. Whether the
      mean is intended is a question for the framework's author.
\item \emph{Start values.} The void fraction is $1-\sum\phi_i$, so start
      values must satisfy $\phi_1+\phi_2\le1$. The framework's example
      input starts both species at $1$ in the seed; for one species I start the
      second at $0$.
\item \emph{The first result set is not in equilibrium.} While the element
      sets up its operators it returns zero stress and tangent on its first
      calls, so the solver converges the first substep on a zero residual. The
      output written for that substep is not an equilibrium state and is
      excluded from every result below.
\item \emph{The front term is switched off, and differs from the paper.}
      An input-reading slip sets the orientation weight of species~1 to zero
      (the weights of both species are read into the same variable, and the
      last one read is zero), so the front-growth term
      of~\eqref{eq:ch5-field} never acts and the biofilm only diffuses. The
      term is also multiplied by $|\Delta\phi|$ where the paper has
      $|\nabla\phi|$. On a test copy with the slip fixed, $|\Delta\phi|$
      makes the explicit update unstable at $\Delta t=0.005$, while
      $|\nabla\phi|$ spreads the seed at the paper's pace but lets
      $\phi$ exceed $1$, and the run breaks down near $T^*=0.22$. The
      overshoot follows from a cell P\'eclet number of about $30$ for this
      advection term on the element's mesh. With the slip fixed and the term
      discretised upwind, $\phi$ stays below $1$, but the biofilm grows
      both towards the nutrient and away from it, far more slowly than in the
      paper. Beyond the seed the nutrient is almost uniform, so
      $\nabla c\approx\mathbf 0$, while the paper's direction
      $\nabla c/|\nabla c|$ remains a unit vector whose orientation is then
      set by round-off. In equation~\eqref{eq:ch5-field} as stated the
      direction is undefined wherever $c$ is uniform, so a finer mesh or
      another scheme does not remove it. I do not regularise it here, because
      that would add a threshold of my own; how the authors treat this limit
      is an open question for them. The checks of \S\ref{sec:ch5-results} do
      not depend on this term (the gradient-free region has no front, and
      both ways of computing $\alpha$ see the same field); how far the
      biofilm spreads does, which is why the example model shows no
      spreading.
\end{enumerate}}

%-----------------------------------------------------------------------
\section{Two species: composition from the point model}
\label{sec:ch5-composition}

% TierA. Klempt 2026; klempt2026_reproduction.py; fritsch2025_cases.py.
\TierA{%
\paragraph{The point model.} Klempt et al.~\cite{Klempt2026ContinuumBacterialGrowth}
describe each species $i$ by its volume fraction $\phi_i$ and its living
fraction $\psi_i$, with the void $\phi_0=1-\sum_i\phi_i$:
\begin{equation}
\begin{aligned}
(\eta_{\phi,i}+\eta_i\psi_i^2)\,\dot\phi_i+\eta_i\phi_i\psi_i\,\dot\psi_i
&= c^*\psi_i\textstyle\sum_j A_{ij}\phi_j\psi_j-\gamma-(\text{barrier}),\\
\eta_i\,(\phi_i\psi_i\,\dot\phi_i+\phi_i^2\,\dot\psi_i)
&= c^*\phi_i\textstyle\sum_j A_{ij}\phi_j\psi_j-\alpha^* b_i\,\psi_i-(\text{barrier}),
\end{aligned}
\label{eq:ch5-point}
\end{equation}
where $A_{ij}$ are the interactions, $b_i$ the sensitivity to an antibiotic of
concentration $\alpha^*$, $c^*$ the nutrient, $\gamma$ the multiplier for
$\sum_i\phi_i+\phi_0=1$, and the barrier terms keep $\phi_i,\psi_i$
in $(0,1)$. My implementation reproduces all fourteen numerical examples of the
paper. Two printed details had to be read differently to do so: Eq.~17 as
printed carries a $+\gamma$ that the figures do not show, and the four-species
interaction matrix of Eq.~19 reproduces the figures without its prefactor
$\tfrac12$. The same model reproduces the figures of Fritsch et
al.~\cite{Fritsch2025BayesianMicrofilms} within $0.013$.}

% TierA. SPLIT_COUPLING.md option (i); composition_reference.py.
\TierA{%
\paragraph{Amount from the field, composition from the point model.} The
element already solves two species fields, so the point model is not needed to
say how much biofilm there is; Klempt et al.\ 2024 define $\phi$ as the
total biofilm fraction, which is $1-\phi_0$ of the point model. I therefore
let the field decide the amount, $\phi_{3D}=\phi_1+\phi_2$, which
drives growth through~\eqref{eq:ch5-eq36}, and let the point model decide only
the composition (the shares $\phi_i/\sum_j\phi_j$) and $\psi_i$. At each
Gauss point and substep its state is scaled to $\sum_i\phi_i=\phi_{3D}$,
advanced, and scaled again; $\phi_i$, $\psi_i$ and $\gamma$ are stored as state
variables. The parameters are those of the paper's Table~1 cases ($c^*=100$,
$\alpha^*=10$, $\eta_i$ per case).

\paragraph{Consistency with the point model.} The scaling changes nothing
when the field's amount is the one the point model would produce on its own.
Driven by that amount, the coupled scheme gives back the composition
of the point model alone: to $1.2\times10^{-5}$ in case~3 and to
$6\times10^{-17}$ in case~6 of the paper, at coupling steps from $0.01$ down to
$0.0025$.%
\footnote{\texttt{tests/test\_composition\_reduces\_to\_point\_model.py}.}
The coupled composition can therefore differ from the paper only where the
field's amount differs from the point model's, and the difference is caused by
that amount and not by the coupling itself. No published work states how the
two models are to be combined, since each prescribes its own evolution of the
amount; a coupling that needs no such choice would require a spatial
multi-species model derived from a single variational principle.

\paragraph{A second scheme.} Klempt et al.\ describe their model as a
material point model and state that $\phi_0$ ``rather serves as a numerical
auxiliary quantity'' and is not to be read as physical voids
\cite{Klempt2026ContinuumBacterialGrowth}. The scaled scheme above uses
$\phi_0=1-\phi_{3D}$ and so gives $\phi_0$ a physical meaning. I
therefore also implemented a second scheme that keeps $\phi_0$ auxiliary:
at each Gauss point the point model runs on its own, unscaled, from the step
in which the field first reaches $\phi_{\min}$, and supplies only the
composition, while the amount remains the field's. Its composition is then
always the paper's trajectory, shifted by the local arrival time of the
biofilm, and needs no cap. In the packed seed at $T^*=1$ it gives
$\phi_1/(\phi_1+\phi_2)=0.614$ in case~3 and $0.014$ in case~6, the values of the point model
alone, against $0.667$ and $0.021$ for the scaled scheme. The two schemes
differ in what they let the composition depend on: the local amount in the
first, the local age of the biofilm in the second.

The composition does not act back on the stress. In Klempt et al.\ 2024 no
material parameter depends on the species; in the earlier spatial model of
Klempt et al.~\cite{Klempt2023NumericalComparison} species identity enters only
through the front growth factor and the nutrient consumption, and I take both
equal for all species. A composition-dependent front factor would be the
natural two-way extension.}

% TierA. SPLIT_COUPLING.md: s sweep, phi_cap, phi_min.
\TierA{%
\paragraph{Assumptions.} Three numbers are not given by either paper, and I
state each as an assumption.
\begin{itemize}
\item \emph{Time link $s$.} Klempt et al.\ 2024 normalise time to
      $T^*\in[0,1]$; the 2026 paper runs $500$ to $1500$ steps of $10^{-4}$ in
      its own unit. No paper links the two. The point model advances
      $s\,\Delta t$ per substep, with $s=0.15$, which maps the 2026 horizon onto
      $T^*=1$.
\item \emph{Cap $\phi_{\rm cap}=0.9$.} Where the field is fully packed,
      $\phi_{3D}\to1$, the point model would be held at a void fraction below
      its own packing limit, and the composition then depends on the coupling
      time step. The point model therefore sees $\min(\phi_{3D},\phi_{\rm cap})$.
\item \emph{Threshold $\phi_{\min}=0.01$.} Started from $\phi=0$ the
      point model creates biomass, because its barrier terms make $\phi=0$ a
      non-equilibrium. It is therefore not run where
      $\phi_{3D}<\phi_{\min}$, and the composition is held there. The local source of~\eqref{eq:ch5-field} raises the void by about
      $k_\alpha T^*\le10^{-3}$, well below the threshold.
\end{itemize}}

% Added 4 Oct 2026: the whole coupled step in one place
% (callsite/phi_mode_exec.inc, composition branch; usermat_py_hook.f;
% coupling/material_server.py).
\TierA{%
\paragraph{The complete step.} Algorithm~\ref{alg:ch5-coupled} collects one
substep of the coupled run. ANSYS solves only the equilibrium
$\operatorname{div}\boldsymbol\sigma=\mathbf 0$ by Newton's method; the
material routine turns the deformation at a Gauss point into a stress after
removing the growth, and, for two species, asks the point model for the new
composition. The composition is stored and does not enter the stress, so it
does not affect the Newton iteration.}

\begin{algorithm}[htbp]
\caption{One substep $t_n\to t_{n+1}$ of the coupled run}
\label{alg:ch5-coupled}
\begin{algorithmic}[1]
\REQUIRE at every Gauss point the converged state of $t_n$: $\alpha_n$ and
  the point-model state $(\phi_i,\psi_i,\gamma)_n$; the element's fields
  $\phi_1,\phi_2,c$ of $t_n$
\ENSURE displacements, stresses and state of $t_{n+1}$
\REPEAT[Newton iteration of ANSYS]
  \STATE compute the deformation gradient $\mathbf F$ at every Gauss point from the current displacements
  \FORALL{Gauss points}
    \STATE $\phi_{3D}\leftarrow\phi_1+\phi_2$ from the element; restore the state of $t_n$
    \STATE $\alpha_{n+1}\leftarrow\alpha_n+k_\alpha\,\phi_{3D}\,\Delta t$ \COMMENT{Eq.~\eqref{eq:ch5-eq36}}
    \STATE $\mathbf F_e\leftarrow\mathbf F\,\mathbf F_g^{-1}$ with $\mathbf F_g=\alpha_{n+1}\mathbf I$
    \STATE stress $\boldsymbol\sigma$ and tangent from $W(\mathbf F_e)$ and $E(\phi)$, Eq.~\eqref{eq:ch5-stiffness}
    \IF{two species and $\phi_{3D}\ge\phi_{\min}$}
      \STATE scale the point-model state to $\sum_i\phi_i=\min(\phi_{3D},\phi_{\rm cap})$
      \IF{this state was already advanced in this substep}
        \STATE take the stored result
      \ELSE
        \STATE send state, parameters, $s\,\Delta t$ and $n_{\rm sub}$ through the bridge (Fortran, C, local socket)
        \STATE the point model, Eq.~\eqref{eq:ch5-point}, takes $n_{\rm sub}$ steps of $s\Delta t/n_{\rm sub}$ and returns $(\phi_i,\psi_i,\gamma)_{n+1}$
      \ENDIF
      \STATE scale back to $\phi_{3D}$; store $\phi_i$, $\psi_i$, $\gamma$ as state variables
    \ENDIF
    \STATE return $\boldsymbol\sigma$ and the tangent to ANSYS
  \ENDFOR
  \STATE assemble the residual and update the displacements
\UNTIL{the residual meets the convergence criterion}
\STATE the element advances its fields $\phi_1,\phi_2,c$ to $t_{n+1}$ (Eq.~\eqref{eq:ch5-field}, its own solver)
\end{algorithmic}
\end{algorithm}

\begin{table}[htbp]
  \centering
  \caption[Sensitivity to the two coupling assumptions]{%
    Share of species~1, $\phi_1/(\phi_1+\phi_2)$, at $T^*=1$ (case~3 of Klempt et al.\ 2026, coexistence;
    case~6, one species wins). Computed with the same scheme the ANSYS runs
    use. ``Front'' is a point the field fills from $0.05$ to $1$.}
  \label{tab:ch5-sensitivity}
  \begin{tabular}{lcccc}
    \hline
    & $s=0.05$ & $s=0.15$ & $s=0.5$ & $s=1$\\
    \hline
    case 3, front & 0.638 & 0.672 & 0.701 & 0.744\\
    case 6, front & 0.230 & 0.014 & 0.014 & 0.014\\
    \hline
    & $\phi_{\rm cap}=0.85$ & $\phi_{\rm cap}=0.9$ & $\phi_{\rm cap}=0.95$ &\\
    \hline
    case 3, packed seed & 0.668 & 0.667 & 0.664 &\\
    case 6, packed seed & 0.023 & 0.021 & 0.014 &\\
    \hline
  \end{tabular}
\end{table}

\TierA{%
Table~\ref{tab:ch5-sensitivity} shows what the assumptions do. The time link
matters: in case~6 the winner only emerges for $s\gtrsim0.15$ and the result
does not change for larger $s$; in case~3 the share keeps rising slowly with $s$,
from $0.672$ at $s=0.15$ to $0.744$ at $s=1$, while the outcome, coexistence,
stays the same. The cap does not: between $0.85$ and $0.9$
the packed seed and the interior agree, so the composition is continuous across
the edge of the packed region. With $\phi_{\rm cap}=0.9$ the result is
converged at a coupling time step of $0.025$ (to $0.0016$ against $0.0125$);
$\phi_{\rm cap}=0.95$ needs $0.0125$.}

%-----------------------------------------------------------------------
\section{Results}
\label{sec:ch5-results}

% TierA. figs_1005.py; ONE_SPECIES_COUPLING.md, re-run with Table 2 values.
\TierA{%
All runs use ANSYS 2022~R2 on the framework's example model (512 elements),
one process, $\Delta t=0.1$, and the growth parameters of Klempt et al.\ 2024,
Table~2: $k_\alpha=10^{-3}$, one species, the second field started at zero.
The model is a 2~mm cube instead of the paper's 20~$\mu$m one. For the
diffusion coefficient of~\eqref{eq:ch5-field} I use
$\beta=0.02$~mm$^2/T^*$, the value $\beta=2$ of Klempt et al.\ 2024, Table~2,
converted to the 2~mm cube. The conversion is an assumption: the paper gives
no length and time scale, and its first author describes $\beta$ as chosen
for sensible results and dependent on the mesh. The framework's example
input, $\beta=10^{-4}$~mm$^2/T^*$, is reported as a sensitivity case. The
remaining coefficients of the field and of the nutrient ($r$, $k$, $d$, $g$)
are those of the example input; with the front term off
(\S\ref{sec:ch5-coupling}), $r$ and $k$ do not act. The runs end at
$T^*=1.1$ with $\Delta t=0.05$ (one species) and $0.025$ (two species);
halving the step changes the mean stress in the seed by $1$ to $2\,\%$. The
two checks against exact solutions below use the example input and
$\Delta t=0.1$, because they need a region without gradients.

\paragraph{Exact solution inside ANSYS.} Where the field has no gradients,
\eqref{eq:ch5-field} reduces to $\dot\phi=k_\alpha\alpha$,
$\dot\alpha=k_\alpha\phi$, with the solution $\phi=\sinh(k_\alpha t)$,
$\alpha=\cosh(k_\alpha t)$. ANSYS follows the exact solution of its own time
stepping to $3\times10^{-11}$ in $\alpha$ (round-off) and $1.1\times10^{-7}$
in $\phi$, about $10^{-4}$ relative, because the element still sees a small
gradient from its neighbours (Figure~\ref{fig:ch5-exact}).

\begin{figure}[htbp]
  \centering
  \includegraphics[width=\linewidth]{\ansysfigdir fig1005_exact.png}
  \caption[Coupled run against the exact solution]{%
    A gradient-free region of the coupled ANSYS run against
    $\phi=\sinh(k_\alpha t)$, $\alpha=\cosh(k_\alpha t)$. Right: the difference to
    the exact solution of the time-stepping scheme.}
  \label{fig:ch5-exact}
\end{figure}

\paragraph{The whole model.} The element's own growth variable and
equation~\eqref{eq:ch5-eq36} give two ways of computing the stress. Over all
512 elements the two von Mises fields have the same pattern, cosine similarity
$1.0000$. The element's variable gives $0.455$ times the stress
of~\eqref{eq:ch5-eq36}, which is $\tfrac12\times\tfrac{10}{11}$: the mean over
both species fields, one of them unused, and a lag of one substep
(Figure~\ref{fig:ch5-whole}).

\begin{figure}[htbp]
  \centering
  \includegraphics[width=\linewidth]{\ansysfigdir fig1005_whole_model.png}
  \caption[Two ways of computing the growth, whole model]{%
    Element von Mises stress from equation~\eqref{eq:ch5-eq36} against the
    element's own growth variable, one point per element.}
  \label{fig:ch5-whole}
\end{figure}}

% TierA. Re-run on IKMHIWI03 with the Klempt 2024 stiffness (commit c0ac081).
\TierA{%
\paragraph{Stresses with the published stiffness.} With the stiffness of
Klempt et al.\ 2024, equation~\eqref{eq:ch5-stiffness} with $E=10$~Pa,
$\nu=0.49$ and a void floor $f=10^{-3}$, $\alpha-1$ in the seeded element
reaches $8.9\times10^{-4}$ at $T^*=1.1$ on $16^3$ elements
(Figure~\ref{fig:ch5-element}). This is less than $k_\alpha\phi t=1.1\times10^{-3}$,
because $\phi$ diffuses out of the seed. Klempt et
al.~\cite{Klempt2024DiffusionDrivenGrowth,Klempt2023NumericalComparison}
report pressure inside the biofilm and a ring of tension at its edge. Grouped
by distance from the seed, the mean stress is $-2.5\times10^{-4}$~Pa in the
seed, $+6.0\times10^{-5}$~Pa in the first layer of elements around it and
$+1.3\times10^{-6}$~Pa in the second; every element of these two layers is in
tension. The von Mises stress of the seeded element is $4.7\times10^{-4}$~Pa
and the largest in the model $7.2\times10^{-4}$~Pa. The element's own growth
variable gives $0.49$ times the growth and the stress of
equation~\eqref{eq:ch5-eq36}, which is $\tfrac12\times1.1/1.125$ for this time
step (see the whole-model check above). With the framework's example
stiffness ($E=1000$~MPa, $\nu=0.3$, blended linearly) the stresses are about
$4\times10^{8}$ times larger. Changing $\nu$ from $0.49$ to $0.45$ or $0.499$
changes the mean stress in the seed by less than $2\,\%$.

The first author of the paper has since suggested that its stiffness unit may
be wrong and that $E=10$~kPa is more realistic. Growth is the only load, so
all stresses in this chapter then scale by $10^3$; the comparisons and ratios
do not change.

\begin{figure}[htbp]
  \centering
  \includegraphics[width=\linewidth]{\ansysfigdir fig_element_beta002.png}
  \caption[The seeded element and its neighbours]{%
    Von Mises stress of the seeded element and of its most loaded neighbour
    (left) and growth $\alpha-1$ of the seeded element (right) over time;
    $16^3$ elements, $\beta=0.02$~mm$^2/T^*$, stiffness of Klempt et al.\ 2024.
    Solid: equation~\eqref{eq:ch5-eq36} in the material routine; dashed: the
    element's own growth variable.}
  \label{fig:ch5-element}
\end{figure}}

% TierA for the Python study (mesh_study_seed.py, locking_check.py);
% ANSYS 16^3 run: RUN_1005_IKMHIWI03.md (5 Oct).
\TierA{%
\paragraph{Mesh and incompressibility.} The example model has $8^3$
elements, and its seed is 32 of them. I solved the same mechanical problem
(seed growth $1.1\times10^{-3}$, the stiffness above) with a linear
hexahedral solver in Python on $8^3$, $16^3$ and $32^3$ elements. The seed
averages of von Mises and mean stress on $8^3$ are about $1.6$ to $1.75$
times those on $16^3$ and $32^3$, which agree with each other; the peaks at
the corners of the seed do not converge, because those corners are
singular. I therefore quote seed averages and not peaks. With the element's
mean-dilatation (B-bar) integration the mean stress hardly changes between
$\nu=0.49$ and $0.499$ ($-3.8$ and $-4.1\times10^{-5}$~Pa on $8^3$), so
there is no volumetric locking; full $2\times2\times2$ integration would
overestimate it by a factor of $3.7$ on $8^3$ and $2.3$ on $16^3$ at
$\nu=0.49$.}

\TierA{%
The ANSYS model uses the same integration (SOLID185, KEYOPT(2) $=0$, B-bar).
In ANSYS I refined the example model to $16^3$ and $24^3$ elements, with every
boundary condition mapped by position ($32^3$ exceeds the memory of the
workstation in the nutrient solve). Table~\ref{tab:ch5-mesh} gives the seed
averages and the first layer around the seed. The mean stress and the ring of
tension agree within $3\,\%$ on the three meshes, and the von Mises stress
changes by $3\,\%$ from $16^3$ to $24^3$.

\begin{table}[htbp]
  \centering
  \caption[Mesh refinement in ANSYS]{Seed averages and the first layer of
    elements around the seed at $T^*=1.1$, $\beta=0.02$~mm$^2/T^*$,
    $\Delta t=0.025$ [Pa].}
  \label{tab:ch5-mesh}
  \begin{tabular}{lccc}
    \toprule
    & $8^3$ & $16^3$ & $24^3$ \\
    \midrule
    mean stress, seed & $-2.52\times10^{-4}$ & $-2.53\times10^{-4}$ & $-2.45\times10^{-4}$ \\
    mean stress, first layer & $+6.01\times10^{-5}$ & $+6.04\times10^{-5}$ & $+5.85\times10^{-5}$ \\
    von Mises, seed & $4.83\times10^{-4}$ & $5.87\times10^{-4}$ & $6.06\times10^{-4}$ \\
    \bottomrule
  \end{tabular}
\end{table}

The diffusion coefficient sets the size of the stresses
(Table~\ref{tab:ch5-beta}): $\phi$ that diffuses out of the seed lowers its
growth near the surface, and the step in $\alpha$ between the inside and the
surface of the seed loads it. From $\beta=0.005$ to $0.1$~mm$^2/T^*$ the mean
stress in the seed falls by a factor of ten. For $\beta\ge0.02$ the diffusion
length $\sqrt{\beta T^*}\ge0.15$~mm is resolved by the meshes and $8^3$ and
$16^3$ agree within $15\,\%$. For smaller $\beta$ the layer is thinner than an
element and the results depend on the mesh. With the example input,
$\beta=10^{-4}$~mm$^2/T^*$, the layer is $0.01$~mm thick: the step in $\alpha$
grows from $0.4$ to $1.8\,\%$ from $8^3$ to $24^3$, and the seed von Mises
stress ($7.5\times10^{-5}$ to $1.5\times10^{-4}$~Pa) does not converge.
With $\beta=0$ the step vanishes and the ratios $16^3/8^3$ are those of the
Python study above.

\begin{table}[htbp]
  \centering
  \caption[Sensitivity to the diffusion coefficient]{Mean stress in the seed
    at $T^*=1.1$ for different $\beta$ [Pa]; $\beta=0.02$ is the value used in
    this chapter.}
  \label{tab:ch5-beta}
  \small
  \begin{tabular}{lccccc}
    \toprule
    $\beta$ [mm$^2/T^*$] & $0.005$ & $0.01$ & $0.02$ & $0.05$ & $0.1$ \\
    \midrule
    $8^3$ & $-2.1\times10^{-4}$ & $-2.7\times10^{-4}$ & $-2.5\times10^{-4}$ & $-1.3\times10^{-4}$ & $-4.3\times10^{-5}$ \\
    $16^3$ & $-3.9\times10^{-4}$ & $-3.4\times10^{-4}$ & $-2.5\times10^{-4}$ & $-1.2\times10^{-4}$ & $-3.8\times10^{-5}$ \\
    \bottomrule
  \end{tabular}
\end{table}}

% TierA. SPLIT_COUPLING.md, phi_cap runs on IKMHIWI03 (judge PASS, diff 0.0).
\TierA{%
\paragraph{Two species.} With the point model coupled, every traced point of
the ANSYS run equals the same scheme run on its own, with a difference of
$0$, and every call to the point model is reproduced bit for bit. In
the packed seed, at $T^*=1$ (example input, $\Delta t=0.1$):
\begin{itemize}
\item case~3 (coexistence): the paper gives $\phi_1=0.60$, $\phi_2=0.40$ for the point model
      alone, and my implementation of the point model alone gives $0.5945$; the
      coupled run gives $0.6665$;
\item case~6 (one species wins): the paper gives $\bar\phi_1\approx0$ with
      $\psi_1=0.05$; the coupled run gives $\phi_1/(\phi_1+\phi_2)=0.0209$ with $\psi_1=0.048$.
\end{itemize}
The type of outcome is kept. The ratio in case~3 shifts because the amount is
now set by the field instead of by the point model's own filling. The second
scheme, run in the same way, gives $\phi_1/(\phi_1+\phi_2)=0.6137$ ($\psi=0.9849/0.9777$) in
case~3 and $0.0138$ ($0.0545/0.9836$) in case~6, the point model's own values,
again with every point equal to the scheme run on its own.

With $\beta=0.02$~mm$^2/T^*$ the field spreads beyond the seed, and the point
model runs wherever $\phi$ exceeds $\phi_{\min}$: in $304$ of $512$ elements
on $8^3$ and $1970$ of $4096$ on $16^3$. The seed keeps the outcome of the
point model alone, $\phi_1/(\phi_1+\phi_2)=0.667$ in case~3 and $0.03$ to
$0.06$ in case~6. Around the seed the composition now varies in space even
without a nutrient gradient, because a point starts its point model only when
the biofilm reaches it: in case~6 the first layer around the seed has a mean
share of $0.33$, the second $0.50$, the start value. This part depends on an
assumption. The papers do not say which composition a point starts from when
the biofilm reaches it; the first author of the 2026 paper suggests the
composition of the neighbouring points. With start values $0.2$, $0.5$ and
$0.8$ for these points the first layer has $0.34$, $0.41$ and $0.61$ (case~6,
consumption $6$), while the seed does not change. Case~6 has two stable
outcomes: a start share below about $0.51$ ends with species~2, above it with
species~1, at every nutrient level (Figure~\ref{fig:ch5-map}).

\begin{figure}[htbp]
  \centering
  \includegraphics[width=\linewidth]{\ansysfigdir fig_composition_map.png}
  \caption[Which species wins in case 6]{%
    Case~6 of Klempt et al.\ 2026 in the coupling scheme of this chapter, one
    packed point (Python): $\phi_1/(\phi_1+\phi_2)$ at $T^*=1$ against the
    local nutrient and the start share, for $s=0.15$ and $0.25$. Black line:
    $0.5$.}
  \label{fig:ch5-map}
\end{figure}

To see what the coupling does on a spreading biofilm, I ran the same scheme
in Python on the reproduced field of Klempt et al.\ 2024, test case~4.1,
with the point model at every grid node (Figure~\ref{fig:ch5-spreading}).
There each point starts its point model when the field first reaches
$\phi_{\min}$, from the seed composition, because the composition is not
transported with the biofilm. Starting it instead from the composition of the
nearest colonised point changes the biomass-weighted mean of
$\phi_1/(\phi_1+\phi_2)$ at $T^*=1$ from $0.656$ to $0.656$ in case~3 and
from $0.026$ to $0.027$ in case~6; the largest pointwise change is $0.093$,
at the points colonised last. The diffusive edge of the field reaches
$\phi_{\min}$ long before the dense front, so every point runs its point
model for most of the time.

\begin{figure}[htbp]
  \centering
  \includegraphics[width=\linewidth]{\ansysfigdir fig_composition_spreading.png}
  \caption[Composition along a spreading front]{%
    Composition along a spreading front, computed in Python with the coupling
    scheme of this chapter on the reproduced field of Klempt et al.\ 2024,
    test case~4.1; Klempt et al.\ 2026, cases~3 and~6.}
  \label{fig:ch5-spreading}
\end{figure}}

%-----------------------------------------------------------------------
\section{Limitations}
\label{sec:ch5-limitations}

\TierA{%
\begin{itemize}
\item \emph{Validation.} Everything in this chapter is verification. No
      measured stress is available to validate against.
\item \emph{One-way composition.} The species composition does not act back on
      the mechanics, because in the formulation followed here no material
      parameter depends on the species.
\item \emph{Assumptions.} The time link $s$, the cap $\phi_{\rm cap}$, the
      converted $\beta$ and the start composition of points the biofilm
      reaches later are not given by the papers. All are reported with their
      sensitivity; the stresses depend on $\beta$, the composition on $s$ and,
      outside the seed, on the start composition.
\item \emph{Pressure term.} The spherical term of
      \S\ref{sec:ch5-verification-limits} is documented, not corrected. It does
      not change the von Mises stress at a given deformation, and at
      $\nu=0.49$ it is $1.3\,\%$ of the constrained pressure.
\item \emph{Incompressibility.} The paper's exact constraint is approximated by
      $\nu=0.49$ with the element's mean-dilatation formulation; the mean
      stress in the seed changes by less than $2\,\%$ between $\nu=0.45$,
      $0.49$ and $0.499$.
\item \emph{Mesh.} With $\beta=0.02$~mm$^2/T^*$ the seed averages agree within
      $3\,\%$ on $8^3$, $16^3$ and $24^3$ elements; with the example input
      $\beta=10^{-4}$ they do not converge. Peak stresses at the seed corners
      are mesh dependent and not reported.
\item \emph{Growth law.} The paper gives two forms of the $\alpha$ law
      (Eq.~36 and Table~1). I use Eq.~36; under Table~1 a full seed would not
      grow.
\item \emph{No growth towards the nutrient.} The element's front term is
      switched off, so the biofilm spreads only by diffusion. The composition
      varies in space through the local nutrient and the time at which the
      biofilm reaches a point, but not through growth towards the nutrient.
\item \emph{Front growth where the nutrient is uniform.} The direction
      $\nabla c/|\nabla c|$ of the front term in Klempt et al.\ 2024 is
      undefined where $\nabla c=\mathbf 0$; the authors' implementation uses
      $\nabla c/(|\nabla c|+\varepsilon)$ with a small $\varepsilon$. Running
      the front term in the element is left to the continuation at Keio.
\item \emph{Magnitude.} With the published growth rate, $\alpha-1$ stays below
      about $10^{-3}$ up to $T^*=1$, and with the published stiffness the
      largest von Mises stress is about $7\times10^{-4}$~Pa ($0.7$~Pa with
      $E=10$~kPa, the value suggested by the paper's first author).
\end{itemize}}

%-----------------------------------------------------------------------
\section{Outlook: towards a two-way coupling}
\label{sec:ch5-outlook}

% TierA. ROADMAP_TWO_WAY.md.
\TierA{%
The coupling of this chapter is one-way: the field drives growth and stress,
and the point model returns a composition that nothing reads back. Four steps
would close the loop, ordered by how well the literature supports them.
\begin{enumerate}
\item \emph{Local nutrient into the point model.} Klempt et
      al.~\cite{Klempt2026ContinuumBacterialGrowth} treat $c^*$ as a given,
      possibly time-dependent, variable; replacing it by the local $c$ of the
      field needs no new physics.
\item \emph{Composition into the spreading.} In the earlier spatial model the
      front factor and the consumption depend on the species
      \cite{Klempt2023NumericalComparison}; $R_s=\sum_i\phi_iR_{s,i}/\sum_j\phi_j$ and
      $g=\sum_i\phi_i g_i/\sum_j\phi_j$ would let the composition shape where the biofilm
      grows. This needs a working front term
      (\S\ref{sec:ch5-coupling}, property~4).
\item \emph{Composition into the stiffness.} $E=\sum_i\phi_iE_i/\sum_j\phi_j$ would let the
      composition reach the stress directly, but species-specific moduli are
      not available \cite{Bol2013BiofilmMechanicalCharacterisation}.
\item \emph{Stress into growth.} The mechanical term of Klempt et al.\ 2024,
      Eq.~30, dropped in their Eq.~34, and stress-dependent growth as in
      \cite{Soleimani2020AnisotropicGrowth} would close the loop from the
      mechanics back to the field.
\end{enumerate}
Steps 1 and 2 are the natural first goal of the continuation at Keio; each
must reduce to the one-way scheme when the species are made equal.}

% Added 4 Oct 2026: ansys_usermat/composition_local_nutrient.py,
% two_way_step2.py, two_way_step3.py; ROADMAP_TWO_WAY.md.
\subsection{First steps, computed in Python}
\label{sec:ch5-outlook-python}

\TierA{%
I computed the first three steps in Python, outside ANSYS, to see which of them
change the results. The papers give no growth rates or moduli per species.
The rates are estimated from the dissipation parameters $\eta_i$ under a stated
assumption, the moduli are sensitivity parameters.

\paragraph{Step 1: local nutrient.} On the example model of this chapter
(Figure~\ref{fig:ch5-local-nutrient}), the nutrient is held at the face
$y=-1$~mm and consumed in the seed; its strength is the Thiele number
$\Lambda^2=gL^2/d$ with $L=1$~mm, scanned because $d$ and $g$ of the example
input are not taken from the paper. Each point of the seed runs the point
model with $c^*=c^*_0\,c$. Case~3 hardly depends on the nutrient level, the
share $\phi_1/(\phi_1+\phi_2)$ staying at $0.63$ to $0.66$. In case~6 the
takeover is slower where the nutrient is low: for $\Lambda=4$ the share ranges
from $0.03$ near the nutrient face to $0.44$ in the interior of the seed. The
composition thus varies in space without any spreading, which the element
cannot show with a constant $c^*$. The values are converged in the grid
($16^3$, $32^3$) and in the coupling step ($0.025$).

The same step then ran in ANSYS. The nutrient that the element solves at each
Gauss point is passed to the point model as $c^*=c^*_0\,\min(c/c_\mathrm{ref},1)$
with $c_\mathrm{ref}=1$, the held value; without this input the new build gives
the earlier results to the last digit. On the $8^3$ example model with the
example input $\beta=10^{-4}$, case~6, with
the point model in all 512 elements and the consumption raised from the
example value $0.5$ to $4$ (an example input, not a paper value), the
nutrient in the seed falls to $0.44$--$0.66$ and the share
$\phi_1/(\phi_1+\phi_2)$ ranges from $0.025$ at the side facing the nutrient
to $0.14$ at the far side (correlation with $c$: $-0.83$); with the example
consumption it is $0.021$ everywhere. Halving the time step changes the mean
share by $2\,\%$. With consumption $6$, the strongest that keeps $c$ positive,
the nutrient falls to $0.15$--$0.48$ and the share ranges from $0.11$ to
$0.48$ (correlation $-0.99$). Case~3 hardly depends on the nutrient
($0.656$--$0.664$ at consumption $4$), as predicted. This is the first composition in the
element that varies in space. Python gives the same: with the consumption set
so that the lowest nutrient in the seed is $0.435$, as in ANSYS, the share
ranges from $0.026$ to $0.139$. This holds whether the consumption is
independent of $c$ (zero order, as in Eq.~35 of Klempt et al.\ 2024 and in the
element) or proportional to $c$ (first order, as in the Python study above). For consumption $6$ the zero-order
Python model gives $0.105$ to $0.481$, again as ANSYS. The share follows the
local nutrient; the order of the consumption changes how $c$ is distributed,
and with it the range once the consumption is strong. With zero-order consumption a stronger consumption ($8$)
drives $c$ below zero, which is outside the valid range of the model.

With $\beta=0.02$~mm$^2/T^*$ the result is the same in the seed: consumption
$6$ gives a share of $0.12$ to $0.46$ on $8^3$ and $0.07$ to $0.47$ on $16^3$,
where the face towards the nutrient is resolved. The consumption of Klempt et
al.\ 2024, Table~2, converted to the cube like $\beta$, is weaker: the nutrient
in the seed stays at $0.88$ to $0.94$ and the share at $0.03$ to $0.06$
($8^3$), close to the run without nutrient. Over all runs the share at the
end depends on the time link $s$, the end time and the local nutrient only
through the product $s\,T^*c$ (Figure~\ref{fig:ch5-clock}): the nutrient acts
as the clock of the point model. In case~6 the takeover by species~2 happens
between $s\,T^*c\approx0.03$ and $0.08$; a spatial pattern appears where the
local values of $s\,T^*c$ straddle this window. The assumed $s$ therefore
decides whether a composition pattern is seen at $T^*=1$, and I report it
with this dependence.

\begin{figure}[htbp]
  \centering
  \includegraphics[width=\linewidth]{\ansysfigdir fig_composition_clock.png}
  \caption[The nutrient as the clock of the point model]{%
    $\phi_1/(\phi_1+\phi_2)$ at the end of every seed element of the $8^3$
    ANSYS runs ($\beta=0.02$~mm$^2/T^*$; consumption $0$ to $6$,
    $s=0.05$ to $1$, $T^*=1$ and $2$) against $s\,T^*c$, with $c$ the local
    nutrient. Lines: one packed point in Python at $T^*=1$. Left case~6, right
    case~3 of Klempt et al.\ 2026.}
  \label{fig:ch5-clock}
\end{figure}

\begin{figure}[htbp]
  \centering
  \includegraphics[width=\linewidth]{\ansysfigdir fig_composition_local_nutrient_thesis.png}
  \caption[Local nutrient in the point model]{%
    Local nutrient in the point model, computed in Python on the example model
    (section through the seed, $32^3$ grid). Top: nutrient $c$, held at $y=-1$~mm
    and consumed in the seed (white outline), for three Thiele numbers $\Lambda$.
    Bottom: share $\phi_1/(\phi_1+\phi_2)$ at $T^*=1$ for case~6 of Klempt et
    al.\ 2026 with $c^*=c^*_0\,c$; grey: no biofilm. Case~3 stays at
    $0.63$--$0.66$ and is not shown.}
  \label{fig:ch5-local-nutrient}
\end{figure}

\paragraph{Step 2: composition into the spreading.} On the reproduced field of
Klempt et al.\ 2024, test case~4.1, the front growth rate becomes
the mean of the species rates weighted by their shares,
$\bar r=\bigl(\phi_1r_1+\phi_2r_2\bigr)/(\phi_1+\phi_2)$, with
$r_{1,2}=(1\pm d)\,r$ and $r$ from Table~\ref{tab:theory_klempt2024}; for $d=0$ the scheme is the one-way one. The
weighting has two published precedents. In the two-species oral biofilm model
of Feng et al.~\cite{Feng2021SymbioticBiofilm}, the biofilm spreads with a
velocity $\mathbf u$ whose divergence is the sum of the species' growth,
$\nabla\cdot\mathbf u=\sum_i\phi_i\mu_i\,(\ldots)$ with growth rates $\mu_i$,
Monod factors $(\ldots)$ and $\sum_i\phi_i=1$ in the biofilm (their Eq.~15, in
the notation used here), which is the same share-weighted rate; there it
drives a potential flow of the whole biofilm, here it scales the front term of
the field. Summing the growth sources $R_{s,i}\,\phi_i\,c$ of Soleimani et
al.~\cite{Soleimani2023CoAggregation} gives the same weighting for a local
reaction. The rate contrast has two estimates. In cases~3 and~6 of Klempt et
al.\ 2026 the second species has twice the dissipation of the first,
$\eta_1=1$, $\eta_2=2$; assuming that the front rate of a species scales with
$1/\eta_i$ gives $r_1/r_2=2$ and $d=1/3$. Feng et al.\ use maximum growth rates
of $3\times10^{-5}$ and $8\times10^{-5}$~s$^{-1}$ for their two species, the
first measured and the second estimated, a ratio of $2.7$ and $d\approx0.45$. With $d=1/3$
the mean $\phi$ at $T^*=1$ rises from $0.704$ to $0.735$ in case~3, where the
faster species is the majority, and falls to $0.648$ in case~6, where the slower
species takes over; locally $\phi$ changes by up to $0.47$
(Figure~\ref{fig:ch5-two-way}). With $d\approx0.45$ from the rates of Feng et al.\ the
two values are $0.746$ and $0.626$, with $d=0.5$ they are $0.750$ and $0.618$. The composition itself hardly changes. Which
species wins then decides how fast the biofilm spreads. The values are
converged in the coupling step. On a grid of half the spacing ($41^3$ nodes,
$0.5~\mu$m) the one-way mean $\phi$ falls from $0.704$ to $0.674$, and the
changes for $d=0.5$ shrink from $+6\,\%$ and $-13\,\%$ to $+5\,\%$ and
$-8\,\%$: the sign and the order of the effect hold, its size depends on the
grid of the reproduced field, for which I keep the paper's $1~\mu$m. In ANSYS
this step needs a working front term.

\begin{figure}[htbp]
  \centering
  \includegraphics[width=\linewidth]{\ansysfigdir fig_two_way_step2_thesis.png}
  \caption[Composition acting on the spreading]{%
    Composition acting on the spreading, computed in Python on the reproduced
    field of Klempt et al.\ 2024, test case~4.1 (mid-plane $z=10~\mu$m),
    with species growth rates $r_{1,2}=(1\pm d)\,r$, $d=1/3$ from $\eta_1=1$, $\eta_2=2$
    of Klempt et al.\ 2026 (assumption: front rate $\propto1/\eta_i$).
    Left and middle: $\phi$ at $T^*=1$ for case~6 without and with the feedback
    (nutrient on the edge at the top right).
    Right: mean $\phi$ over time.}
  \label{fig:ch5-two-way}
\end{figure}

\paragraph{Step 3: composition into the stiffness.} With the seed composition
of step~1 and a ratio of the species moduli up to $5$ at the same mean modulus,
the average stresses in the seed change by at most $0.2\,\%$. The seed is about
a thousand times stiffer than the void around it, so the soft surroundings set
its stress. This step matters only where the biofilm itself carries the load,
such as a contiguous biofilm on a substrate.

\paragraph{Step 4: stress into growth.} In the example model the mean stress
in the seed is about $10^{-5}$ times the shear modulus, so a stress-dependent
growth law would change little with the published parameters, while the
mechanical term of Klempt et al.\ 2024, Eq.~30, kept as written would stop all
growth (Appendix~\ref{app:klempt2024}). This step needs a different growth law
and realistic stress levels.

Of the four steps, step~2 changes the results most, and step~1 already runs in
the element.}

